\documentclass[12pt]{article}
\usepackage[utf8]{inputenc}
\usepackage[spanish]{babel}
\usepackage{graphicx}
\usepackage{geometry}
\usepackage{fancyhdr}
\usepackage{lastpage}
\usepackage{booktabs} 
\usepackage{caption} % Necesario para usar \captionof

\geometry{
    top=4cm, 
    bottom=3cm, 
    left=2cm, 
    right=2cm, 
    headheight=2.5cm, 
    headsep=0.5cm
}

\begin{document}

% ==========================================
%               CARÁTULA
% ==========================================

\begin{titlepage}
    \thispagestyle{empty} 
    \centering
    \vspace*{2cm}
    {\Huge \textbf{(( titulo ))}} \\[0.5cm]
    {\Large (( subtitulo ))} \\[1cm]
    \rule{\linewidth}{1pt} \\[1.5cm] 
    
    \begin{flushleft}
        \large
        \textbf{Modelo:} (( modelo )) \\
        \textbf{Tipo:} (( tipo )) \\
        \textbf{Fecha:} (( fecha ))
    \end{flushleft}
\end{titlepage}

% --- PÁGINA DE CONTENIDO ---
\newpage
\pagestyle{fancy}
\fancyhf{}
\fancyhead[L]{\includegraphics[height=1.5cm]{(( logo_path ))}} 
\fancyfoot[R]{Página \thepage\ de \pageref{LastPage}}
\renewcommand{\headrulewidth}{0.4pt}

% ==========================================
%           SECCIÓN DE SALIDA
% ==========================================
\BLOCK{ if show_salida }
\section*{Análisis de Salida}

\begin{figure}[h!]
    \centering
    % Gráfico Principal: IC vs VCE (Más grande o a la izquierda)
    \begin{minipage}{0.48\textwidth}
        \centering
        \includegraphics[width=\linewidth]{(( plot_salida_ic ))}
        \caption{Familia de Curvas: $I_C$ vs $V_{CE}$}
    \end{minipage}
    \hfill
    % Gráfico Secundario: IB vs VCE (Para verificar cte)
    \BLOCK{ if plot_salida_ib }
    \begin{minipage}{0.48\textwidth}
        \centering
        \includegraphics[width=\linewidth]{(( plot_salida_ib ))}
        \caption{Control: $I_B$ vs $V_{CE}$}
    \end{minipage}
    \BLOCK{ endif }
\end{figure}

% TABLA DE DATOS (Inmediatamente después de las figuras)
\begin{center}
    {\small (( tabla_salida )) }
    \captionof{table}{Muestra de datos medidos en zona activa (Inicio, Medio, Fin de cada curva).}
\end{center}

\BLOCK{ endif }

% ==========================================
%           SECCIÓN DE ENTRADA
% ==========================================
\BLOCK{ if show_entrada }
\newpage
\section*{Análisis de Entrada}

% Gráficos de Tensión y Corriente
\begin{figure}[h!]
    \centering
    % Primera fila de gráficos
    \begin{minipage}{0.48\textwidth}
        \centering
        \includegraphics[width=\linewidth]{(( plot_entrada_ib ))}
        \caption{$I_B$ vs $V_{BE}$}
    \end{minipage}
    \hfill
    \begin{minipage}{0.48\textwidth}
        \centering
        \includegraphics[width=\linewidth]{(( plot_entrada_vbe_ic ))}
        \caption{$V_{BE}$ vs $I_C$}
    \end{minipage}

    \vspace{0.3cm}

    % Gráfico de Transferencia solo (o podrias ponerlo arriba si entra)
    \begin{minipage}{0.6\textwidth}
        \centering
        \includegraphics[width=\linewidth]{(( plot_entrada_trans ))}
        \caption{$I_C$ vs $V_{BE}$}
    \end{minipage}
\end{figure}

% TABLA ESPECÍFICA PARA ESTOS GRÁFICOS
\begin{center}
    {\small (( tabla_entrada_vi )) }
    \captionof{table}{Datos medidos de Tensión y Corriente (Muestra)}
\end{center}

\clearpage % Para separar HFE

% Análisis de Ganancia
\section*{Análisis de Ganancia de Corriente DC}

\begin{figure}[h!]
    \centering
    \includegraphics[width=0.85\textwidth]{(( plot_entrada_hfe ))}
    \caption{Evolución del parámetro $h_{FE}$ según corriente de colector.}
\end{figure}

% TABLA ESPECÍFICA PARA HFE
\begin{center}
    \vspace{-0.5cm}
    {\small (( tabla_entrada_hfe )) }
    \captionof{table}{Valores calculados de $h_{FE}$ e incertidumbre asociada.}
\end{center}

\BLOCK{ endif }

% ==========================================
%           HOJA DE FÓRMULAS
% ==========================================

\clearpage

\section*{Sección de fórmulas: Funciones de medición e incertidumbres}
\noindent En esta sección se resumen las expresiones matemáticas fundamentales utilizadas para la caracterización del transistor y el cálculo de sus incertidumbres asociadas.

\subsection*{Funciones de Medición}
Se asume el emisor conectado directamente a masa ($V_e = 0V$). Las tensiones y corrientes se determinan mediante:

\begin{itemize}
    \item \textbf{Tensión de Colector ($V_c$):} 
    \[ V_c = \frac{C \cdot V_{ref}}{2^{N_{ADS}}-1} \]
    \item \textbf{Tensión de Base ($V_b$):} 
    \[ V_b = \frac{C \cdot V_{ref}}{2^{N_{ADS}}-1} \]
    \item \textbf{Corriente de Base ($I_b$):} 
    \[ I_b = \frac{V_2 - V_1}{R_{shuntB}} \] 
    \item \textbf{Corriente de Colector ($I_c$):} 
    \[ I_c = c \cdot \frac{V_{ref}}{2^{N_{INA}}} \cdot \frac{1}{R_{shunt_{INA}}} \] 
    \item \textbf{Ganancia de Corriente ($HFE$):} 
    \[ HFE = \frac{I_c(Q)}{I_b(Q)} \] 
\end{itemize}

\subsection*{Fórmula de Incertidumbres}

\begin{itemize}
    \item \textbf{Incertidumbre de la Tensión de Colector ($u_r(V_c)$):} 
    \[ u_r(V_c) = \sqrt{u_r(c_{ADC})^2 + u_r(V_{ref})^2} \]
    \item \textbf{Incertidumbre de la Corriente de Colector ($u_c(I_c)$):} 
    \[ u_c(I_c) = \sqrt{u_r(c_{INA})^2 + u_r(R_{INA})^2} \cdot \overline{I_c} \]
    \item \textbf{Incertidumbre de la Corriente de Base ($u_c(I_b)$):} 
    \[ u_c(I_b) = \sqrt{\left(\frac{u(V_2)}{R_{shuntB}}\right)^2 + \left(\frac{-u(V_1)}{R_{shuntB}}\right)^2 + \left(\frac{V_2-V_1}{R_{shuntB}^2} u(R_{shuntB})\right)^2} \]
\end{itemize}

% ==========================================
%           PÁGINA DE FIRMAS
% ==========================================
\newpage
\noindent \textit{Nota: Todas las tablas muestran el punto inicial, medio y final de cada barrido realizado.}

\vspace{0.5cm}
\noindent Los ensayos fueron realizados bajo una temperatura de \textbf{(( temperatura )) °C} y una humedad de \textbf{(( humedad )) \%}.
\section*{Personal Técnico Interviniente}
La presente medición y análisis ha sido avalada por el siguiente equipo técnico:

\vspace{1.5cm} 

\noindent
\begin{tabular}{cc}
\BLOCK{ for tecnico in tecnicos }
    \begin{minipage}{0.48\textwidth}
        \centering
        \vspace{1.5cm}
        \rule{6cm}{0.4pt} \\ % Línea de firma
        \textbf{(( tecnico.nombre ))} \\
        (( tecnico.contacto )) \\
        \textit{Firma Digital}
    \end{minipage} 
    \BLOCK{ if loop.index is divisibleby 2 } \\[1.5cm] \BLOCK{ else } \hspace{0.2cm} \BLOCK{ endif }
\BLOCK{ endfor }
\end{tabular}

\end{document}